\subsection{The exact local image of cyclic auxiliary repair}

\begin{statement}{A.122}{Proved}
{A complete mod-three criterion for the relative norm module}
Let \(a\in\Q^\times\) satisfy \(v_3(a)=1\), let \(n\ge1\),
and retain the relative torus and its actual module from A.115:
\[
 \begin{gathered}
 N=3^n,\quad d=N/3,\quad K_j=\Q(a^{1/3^j}),\quad K_0=\Q,\\
 T_n=\ker(\operatorname{Res}_{K_n/\Q}\mathbf G_m
                 \longrightarrow\operatorname{Res}_{K_{n-1}/\Q}\mathbf G_m),
 \qquad C=T_n[3]=Z^d\F_3[Z]/Z^N\otimes\chi,\\
 L=K_n(\mu_N),\qquad G=\operatorname{Gal}(L/\Q).
 \end{gathered}
 \tag{A.122.1}
\]
Let \(E/\Q\) be any finite Galois extension containing \(L\).
Let \(S\) be finite and contain the real place, three, and
every prime ramified in \(E\).
Let \(\mathcal P\) consist of the primes outside \(S\)
splitting completely in \(E\).
Define
\[
 \begin{gathered}
 \mathcal H=\bigcup_{\substack{P\subset\mathcal P\\P\ {\rm finite}}}
                         H^1(G_{\Q,S\cup P},C),\\
 \mathcal H_{\rm cyc}=\{\gamma\in\mathcal H:
       \operatorname{res}_p\gamma\text{ has cyclic image at every }
                                      p\notin S\text{ where it ramifies}\}.
 \end{gathered}
 \tag{A.122.2}
\]
At those ramified places the coefficient action is trivial,
so cyclic image is the literal image of the local
homomorphism into \(C\).
For the basis \(B_j\), \(0\le j<n\), of A.115, put
\[
 Q_{j,w}(\eta)=
       \operatorname{inv}_w\bigl(\tfrac12(B_j)_*(\eta\smile\eta)\bigr).
 \tag{A.122.3}
\]
Then the following equality of actual local images holds:
\[
 \operatorname{loc}_S(\mathcal H_{\rm cyc})=
 \left\{(\eta_w)\in\operatorname{loc}_S(\mathcal H):
             \sum_{w\in S}Q_{j,w}(\eta_w)=0
                                \text{ for every }0\le j<n\right\}.
 \tag{A.122.4}
\]

In particular, prescribe arbitrary classes at finitely
many primes splitting \(E\), and prescribe zero at every
other controlled place, including three.
They have a lift with cyclic new ramification at primes
splitting \(E\) exactly when their total values in
(A.122.3) vanish for every \(j\).
Purely ramified classes at the prescribed split primes
always satisfy that condition.

If the prescribed classes in \(S\) are Kummer classes
of actual local points of \(T_n\), the class supplied
by (A.122.4) is the Kummer class of an actual
\[
             U\in K_n^\times,\qquad
                         N_{K_n/K_{n-1}}U=1 .
 \tag{A.122.5}
\]
The theorem concerns precision three.
It does not identify these finite local tests with
the higher-power norm conditions at an already
allocated coefficient prime.
\end{statement}

\begin{proof}
\emph{The finite-module support comparison.}
For any finite module \(A\) split by a finite Galois
field \(F/\Q\), with all its ramification and coefficient
primes in a finite set \(S_0\), the argument of
A.55 and A.117 gives
\[
 \operatorname{loc}_{S_0}
 H^1_{S_0\cup{\rm cs}(F)}(\Q,A)
   =\left(\operatorname{res}_{S_0}
        H^1(\operatorname{Gal}(F/\Q),A^\vee)\right)^\perp .
 \tag{A.122.6}
\]
The pairing is the sum of the actual local cup-product
invariants.  Finite Poitou--Tate duality identifies the
annihilator for a finite set of new primes with the
fixed-\(S_0\) dual classes locally zero at those primes.
That dual group is finite; its decreasing kernels
stabilize.  Chebotarev in the finite cocycle field
identifies their intersection with the inflated group
in (A.122.6).
See \cite[I, Theorem 4.10(b) and Corollary 4.15]{BIR_MilneADT}.
This also proves (A.122.6) after deleting any fixed
finite collection of permitted test primes.

Necessity in (A.122.4) follows from Brauer reciprocity.
At a new split prime a cyclic class takes values in
one coefficient line, on which every alternating
\(B_j\) vanishes.  At an unramified place the cup
inflates from a procyclic quotient with zero
degree-two torsion cohomology.  Thus all contributions
outside \(S\) are zero.

\emph{A generator at one additional prime.}
Take \(\eta\) in the right side of (A.122.4) and
choose an actual \(\gamma^{(0)}\in\mathcal H\)
localizing to it.
For a representing cocycle let \(F_{\gamma^{(0)}}\)
be the finite Galois field cut out by
\[
 G_\Q\longrightarrow C\rtimes G,\qquad
                  g\longmapsto(\gamma^{(0)}(g),\bar g).
 \tag{A.122.7}
\]
It contains \(L\).
Choose a prime \(q_0\) splitting completely in
\(E F_{\gamma^{(0)}}\), outside every prime so far
used, and set \(t_0=Z^d\in C\).
Its translates span \(C\), since the translation
group acts by the one Jordan block \(1+Z\).

Apply (A.122.6) with the larger splitting field
\(E F_{\gamma^{(0)}}\).
Prescribe a pure tame class with inertia \(t_0\)
and Frobenius zero at \(q_0\), and zero at \(S\)
and at every prime ramified in \(E F_{\gamma^{(0)}}\).
Every dual class in (A.122.6) is locally zero at
\(q_0\), so these prescriptions are in the stated
annihilator and have a global solution.
Add that solution to \(\gamma^{(0)}\), and call
the result \(\gamma\).
It has the same old local tuple \(\eta\), has
pure tame class \(t_0\) at \(q_0\), and all its
other new ramification still splits \(E\).

Let \(F_\gamma\) be defined by (A.122.7) with this
new cocycle, and put \(\Omega=E F_\gamma\).
The restriction of \(\gamma\) to \(G_E\) is a
homomorphism to \(C\).
Its image contains \(t_0\), because \(q_0\) splits
\(E\), and is stable under \(G\), by the cocycle
identity applied to conjugation.
It is therefore all of \(C\).
This proves, through the actual restriction maps,
\[
 F_\gamma\cap E=L,\qquad
 \operatorname{Gal}(\Omega/L)
       =\operatorname{Gal}(E/L)\times C.
 \tag{A.122.8}
\]
In particular the second equality is a direct
product, not just an equality of orders.

\emph{The local corestriction problem.}
Let \(T\) contain \(S\), \(q_0\), and every prime
ramified in \(\Omega\).
Let \(\mathcal R\) be the primes at which \(\gamma\)
ramifies outside \(S\cup\{q_0\}\).
All members of \(\mathcal R\) split \(E\), hence \(L\).
At each such prime choose normalized tame and
unramified coordinates and write
\[
                    \gamma_p=(t_p,c_p).
 \tag{A.122.9}
\]
Choose their orientation so that the local invariant
of \(\tfrac12 B_*(\gamma_p\smile\gamma_p)\), with any
equivariant alternating form \(B\) in place of \(B_*\),
is \(B(t_p,c_p)/3\).
This fixes the scalar identification in the
one-dimensional uniformizer--unit pairing.

We seek
\[
 y\in H^1(\Q_{{\rm cs}(\Omega)\cup T}/L,C)
 \tag{A.122.10}
\]
which is unramified above \(T\), is zero above \(S\),
and whose corestriction is the unramified class
\(-c_p\) at \(p\in\mathcal R\), and zero at \(q_0\).
The additional primes of \(T\) have zero local
prescription.

For each split \(p\in\mathcal R\cup\{q_0\}\),
choose a prime of \(L\) above \(p\) and index
its conjugates by \(\sigma\in G\).
Transport a local unramified class back to this
prime using \(\sigma^{-1}\), including its action
on coefficients.  Write its Frobenius value
as \(a_{p,\sigma}\in C\).
Local corestriction is the literal sum
\[
 \sum_{\sigma\in G}a_{p,\sigma}=-c_p
                    \quad(p\in\mathcal R),\qquad
 \sum_{\sigma\in G}a_{q_0,\sigma}=0 .
 \tag{A.122.11}
\]
These equations define a nonempty finite-dimensional
affine space of local prescriptions.

Over \(L\) both \(C\) and \(C^\vee\) are trivial.
The dual finite-group classes controlling
(A.122.10) are therefore, by (A.122.8),
\[
 \operatorname{Hom}(\operatorname{Gal}(E/L),C^\vee)
                  \ \oplus\ \operatorname{Hom}_{\F_3}(C,C^\vee).
 \tag{A.122.12}
\]
This is the same restricted-support comparison as
(A.122.6), now over \(L\).
Its allowed primes lie over rational primes
splitting \(\Omega\).
They still detect every element of the kernel
over \(\Omega\): take a finite Galois closure
over \(\Q\) of the dual cocycle field and apply
Chebotarev to an element fixing \(\Omega\).
Thus (A.122.12) uses exactly the permitted primes.

The first summand contributes zero to every
prescription under consideration.
At the primes where a prescription can be nonzero,
\(E\) is completely split; at every other prime
the prescription is zero.
For \(f:C\to C^\vee\) in the second summand set
\[
 B_f(v,w)=\operatorname{ev}(v,f(w)),\qquad
 B_f^\sigma(v,w)=\sigma^{-1}B_f(\sigma v,\sigma w).
 \tag{A.122.13}
\]
After identifying \(\mu_3\) with \(\F_3\), the
last expression is
\(\chi(\sigma)^{-1}B_f(\sigma v,\sigma w)\).
Naturality of local invariants and corestriction,
with the transported coordinates above, makes
the entire dual obstruction
\[
       -\sum_{p\in\mathcal R\cup\{q_0\}}
                     \sum_{\sigma\in G}
                           B_f^\sigma(a_{p,\sigma},t_p).
 \tag{A.122.14}
\]
The sign is the local pairing
\((t,c),(u,v)\mapsto t\cdot v-c\cdot u\):
the first class here is unramified.
The formula retains the cyclotomic action on
the value of the pairing.

\emph{The affine annihilator calculation.}
Let \(\mathcal V_0\) be the space of variations
of (A.122.11), whose sum at each prime is zero.
The functional (A.122.14) vanishes on
\(\mathcal V_0\) exactly when \(B_f\) is
\(G\)-equivariant.
One direction follows from the zero sums.
For the other use only the variations at \(q_0\).
They imply
\[
 B_f^\sigma(v,t_0)=B_f(v,t_0)
                      \quad(v\in C,\ \sigma\in G).
 \tag{A.122.15}
\]
Equivalently
\(B_f(w,\rho t_0)=\chi(\rho)B_f(\rho^{-1}w,t_0)\).
Applying this with \(\rho=\sigma\tau\) and
\(\rho=\tau\) gives
\[
 B_f^\sigma(v,\tau t_0)=B_f(v,\tau t_0)
                         \quad(v\in C,\ \sigma,\tau\in G).
 \tag{A.122.16}
\]
The orbit of \(t_0\) spans \(C\), so this is
equivariance on the whole module.

Every equivariant bilinear form on this module
is alternating.  Indeed \(B_*\) of A.115
identifies it with an endomorphism commuting
with \(G\).  The translation centralizer is
the polynomial algebra in \(Z\), and the
arithmetic involution acts on this algebra
by the adjoint substitution
\(Z\mapsto-Z/(1+Z)\).
Commutation makes the endomorphism self-adjoint.
In odd characteristic its pairing through
\(B_*\) is alternating.
A.115 thus identifies the entire annihilator
in (A.122.15) with
\(\langle B_0,\ldots,B_{n-1}\rangle\).

For an equivariant \(B\), (A.122.14) has the
same value on every affine prescription:
\[
                -\sum_{p\in\mathcal R}B(t_p,c_p).
 \tag{A.122.17}
\]
Global reciprocity for
\(\tfrac12 B_*(\gamma\smile\gamma)\), with \(B\)
in place of the displayed pairing, makes
(A.122.17) zero.
The old sum on \(S\) is zero by (A.122.4),
the class at \(q_0\) is pure tame, and every
other contribution is unramified.

Here is the exact affine existence argument.
Choose one solution \(a^{(0)}\) of (A.122.11).
Let \(D\) send a local prescription to the
functional (A.122.14) on
\(\operatorname{Hom}(C,C^\vee)\).
The annihilator of \(D(\mathcal V_0)\) has
just been identified with the equivariant
forms.  Equation (A.122.17) says that
\(D(a^{(0)})\) vanishes on this annihilator.
Finite-dimensional perfect duality places
\(-D(a^{(0)})\) in \(D(\mathcal V_0)\).
Adding the corresponding variation gives
a prescription annihilating all of (A.122.12).
The restricted-support Poitou--Tate comparison
produces the actual class \(y\) of (A.122.10)
with these local restrictions.

\emph{The cyclic correction.}
Apply Schmidt--Wingberg's cohomological theorem
\cite[Theorem 13, pp.~14--19]{RootSchmidtWingberg1998}
to \(\Omega/L/\Q\), the module \(C\), and \(y\).
The field \(L\) contains \(\mu_3\).
The set \(T\) contains all ramification of
\(\Omega/\Q\), three, and the infinite place.
The class \(y\) is unramified above \(T\)
and zero above
\(\operatorname{Ram}(L/\Q)\cup\{3,\infty\}\),
because these places belong to \(S\).
Every hypothesis of the theorem is therefore
satisfied.
It gives a global \(x\), ramified only in
\(T\cup{\rm cs}(\Omega)\), with local values
on \(T\) equal to those of
\(\operatorname{cor}_{L/\Q}y\), and cyclic
local class outside \(T\).

Set \(\gamma'=\gamma+x\).
It has the original restrictions \(\eta\)
on \(S\) and remains pure tame at \(q_0\).
At each \(p\in\mathcal R\) its coordinates
are exactly \((t_p,0)\).
At a newly ramified prime outside \(T\),
the prime splits \(\Omega\); hence \(\gamma\)
is locally zero there by the definition of
its finite cocycle field.
Consequently \(\gamma'_p=x_p\) is cyclic.
The coefficient module at this prime is
trivial, so the cited cyclicity means
cyclic image of the local homomorphism.
All these primes split \(E\).
This proves sufficiency and (A.122.4).

\emph{Split old plans and actual norm-one elements.}
If the old nonzero prescriptions occur only
at split-\(E\) primes, every finite dual class
in (A.122.6) is zero there.
Zero at the other controlled places therefore
puts every such tuple in the unrestricted
image.  Equation (A.122.4) leaves precisely
the quadratic criterion.
Pure tame classes have zero values for all
alternating coefficient forms.

Suppose finally that the old tuple consists
of local torus Kummer classes.
At new split primes the torus is split.
At every other uncontrolled prime the
produced class is unramified and is a local
Kummer class of an integral torus point,
as in A.117.
The relative \(S_3\) extension has full
decomposition at three.
Its Hasse norm principle and the explicit
restriction-of-scalars comparison of A.117
give \(\operatorname{Sha}^1(\Q,T_n)=0\).
The torus Kummer sequence then lifts
\(\gamma'\) to an actual element (A.122.5).
\end{proof}

The retained certificate
\path{code/relative_cyclic_repair_certificate.py}
checks the transported action, the annihilator
of the zero-corestriction distributions, and
the resulting finite affine systems.
The theorem supplies cyclic local images at
precision three.  At a new coefficient prime
with \(s=v_3(p-1)>1\), a later norm in the
same radical tower can still require a symbol
of order \(3^s\); the displayed criterion does
not identify that symbol with its reduction.
